\section{Tampilan Perangkat Lunak}

Pada bab ini akan membahas antarmuka pengguna (\textit{user interface}) dari aplikasi {Data Capture}. Aplikasi ini dirancang untuk memfasilitasi proses pencatatan data penyewa, peninjauan data yang telah tersimpan, serta pembaruan informasi secara dinamis. Hak akses aplikasi dibedakan menjadi dua peran (\textit{role}) utama, yaitu Pengelola dan Pelaku Komersil. Rincian fungsionalitas dan tampilan dari setiap halaman aplikasi dipaparkan sebagai berikut.

\subsection{Halaman \textit{Login}}
Halaman {login} berfungsi sebagai antarmuka autentikasi awal bagi seluruh pengguna. Pada tahap ini, pengguna diwajibkan untuk memasukkan {e-mail} dan kata sandi (\textit{password}).

\begin{figure}[H]
    \centering
    \includegraphics[width=0.85\textwidth]{Gambar/login (sementara).png}
    \caption{Tampilan Halaman \textit{Login}}
    \label{fig:ui_login}
\end{figure}

\subsection{Halaman Daftar}
Halaman pendaftaran (\textit{register}) disediakan secara khusus bagi pengguna dengan peran Pelaku Komersil yang belum memiliki akun, sehingga mereka dapat mendaftarkan diri secara mandiri. Sementara itu, akun dengan peran Pengelola tidak dapat mendaftar melalui halaman ini karena akun Pengelola akan didaftarkan langsung oleh administrator.

\begin{figure}[H]
    \centering
    \includegraphics[width=0.85\textwidth]{Gambar/daftar (sementara).png}
    \caption{Tampilan Halaman Pendaftaran Akun}
    \label{fig:ui_daftar}
\end{figure}

\subsection{Halaman \textit{Dashboard}}
Halaman {dashboard} merupakan antarmuka utama yang ditampilkan setelah pengguna berhasil masuk (\textit{login}) ke dalam aplikasi. Tampilan dan informasi pada halaman ini disesuaikan dengan hak akses dari masing-masing peran pengguna.

Bagi pengguna dengan peran Pengelola, halaman {dashboard} menyajikan data statistik penyewa secara keseluruhan, mencangkup seluruh aktivitas transaksi yang dilakukan oleh semua Pelaku Komersil. Selain itu, Pengelola juga memiliki wewenang khusus untuk mengunggah atau melakukan perubahan terhadap Laporan Pelanggaran mencangkup foto dan deskripsi pelanggarannya. Laporan yang telah diunggah tersebut kemudian dapat dilihat oleh seluruh pengguna aplikasi.

Sementara itu, bagi pengguna dengan peran Pelaku Komersil, data statistik yang ditampilkan pada {dashboard} dibatasi hanya pada rekapitulasi data penyewaan dari akun miliknya sendiri. Terkait fitur Laporan Pelanggaran, Pelaku Komersil hanya memiliki hak akses untuk melihat laporan yang telah dipublikasikan oleh Pengelola dan tidak memiliki kewenangan untuk mengunggah laporan pelanggaran.

\begin{figure}[H]
    \centering
    \includegraphics[width=0.85\textwidth]{Gambar/dashboard sesudah.png}
    \caption{Tampilan Halaman {Dashboard}}
    \label{fig:ui_dashboard}
\end{figure}


\subsection{Halaman Input}
Halaman \textit{input} didesain untuk mengakomodasi proses rekam data penyewa baru ke dalam basis data aplikasi. Struktur antarmuka pada halaman ini bersifat seragam untuk kedua peran pengguna. Fokus fungsionalitasnya ditujukan untuk mengumpulkan formulir yang mencakup informasi profil identitas penyewa beserta rincian detail transaksi sewanya.

\begin{figure}[H]
    \centering
    \includegraphics[width=0.85\textwidth]{Gambar/input sesudah.png}
    \caption{Tampilan Halaman Input Pengelola}
    \label{fig:ui_input_Pengelola}
\end{figure}

\begin{figure}[H]
    \centering
    \includegraphics[width=0.85\textwidth]{Gambar/input sesudah.png}
    \caption{Tampilan Halaman Input Pelaku Komersil}
    \label{fig:ui_input_Pelaku_Komersil}
\end{figure}

\subsection{Halaman Lihat / Edit Data}
Halaman ini memungkinkan pengguna meninjau ulang riwayat data penyewa yang telah tersimpan, sekaligus melakukan pencatatan informasi pelengkap seperti pencatatan waktu keluar (\textit{check-out}) atau modifikasi metode pembayaran dan jenis sewa. Berikut antarmuka yang dibedakan berdasarkan peran yang diembannya.

\begin{figure}[H]
    \centering
    \includegraphics[width=0.85\textwidth]{Gambar/lihat edit data sesudah.png}
    \caption{Tampilan Halaman Lihat / Edit Data Pengelola}
    \label{fig:ui_lihat_edit_data_Pengelola}
\end{figure}

\begin{figure}[H]
    \centering
    \includegraphics[width=0.85\textwidth]{Gambar/lihat edit data sesudah.png}
    \caption{Tampilan Halaman Lihat / Edit Data Pelaku Komersil}
    \label{fig:ui_lihat_edit_data_Pelaku_Komersil}
\end{figure}

Berikut adalah rincian fitur pada halaman Lihat / Edit Data yang dibagi berdasarkan masing-masing peran.
\begin{itemize}
    \item \textbf{Pengelola:} memiliki kewenangan akses secara menyeluruh. Pengelola berhak meninjau semua rekam data penyewaan yang terdaftar di dalam aplikasi. Melalui halaman ini, Pengelola dapat melengkapi jam \textit{check-out}, memodifikasi metode pembayaran, mengubah jenis sewa, serta melakukan \textit{checklist} pada bagian komentar apabila ditemukan indikasi pelanggaran. Tanda \textit{checklist} ini dapat dibiarkan kosong apabila tidak ditemukan indikasi pelanggaran. Selain itu, Pengelola juga dapat mengunduh (\textit{download}) seluruh data penyewa beserta detail riwayat aktivitasnya (\textit{log}), memverifikasi keberadaan data ganda (\textit{duplicate}), serta menyaring tabel (\textit{filter}) agar hanya menampilkan data penyewa yang terindikasi memiliki pelanggaran. Selain itu, halaman ini juga menyediakan fungsionalitas untuk penyegaran data tabel (\textit{refresh}) dan fitur pencarian informasi spesifik (\textit{search}).
    
    \item \textbf{Pelaku Komersil:} memiliki ruang lingkup kewenangan yang dibatasi. Pengguna dengan peran ini hanya diizinkan untuk mengakses dan meninjau data penyewa yang didaftarkan oleh akunnya sendiri. Pelaku Komersil diberikan hak pembaruan yang sama—seperti menginput waktu \textit{check-out}, memperbarui metode pembayaran, memodifikasi jenis sewa, dan menandai \textit{checklist} komentar pelanggaran (dibiarkan kosong jika tidak terdapat pelanggaran)—namun tindakan tersebut hanya berlaku secara spesifik terhadap data penyewa kepemilikannya. Otoritas pengunduhan data beserta rincian \textit{log}-nya juga dibatasi ketat pada rekam data yang berada di bawah otoritas operasionalnya. Meski demikian, fitur navigasi dasar seperti penyegaran tabel (\textit{refresh}) dan fungsi pencarian data (\textit{search}) tetap dapat diakses secara penuh.
\end{itemize}

\subsection{Halaman Manajemen Unit}
Halaman Manajemen Unit difungsikan sebagai pusat pengelolaan data untuk menambahkan atau mengurangi daftar unit penyewaan yang dimiliki oleh setiap Pelaku Komersil. Halaman ini bersifat eksklusif hanya untuk pengguna dengan peran Pengelola.

\begin{figure}[H]
    \centering
    \includegraphics[width=0.85\textwidth]{Gambar/manajemen unit sesudah.png}
    \caption{Tampilan Halaman Manajemen Unit}
    \label{fig:ui_manajemen_unit}
\end{figure}